\subsection{Symmetric functions in the roots of a polynomial}

Consider a monic polynomial of degree \(n\) , with coefficients in \(A\) :

\[
f = \mathbf {X} ^ {n} + a _ {1} \mathbf {X} ^ {n - 1} + \dots + a _ {n - 1} \mathbf {X} + a _ {n}.
\]

We define an associative, commutative and unital A-algebra \(\mathbf{E}_f\) with the generators \(x_{1},\ldots ,x_{n}\) and relations

\[
\sum_ {i _ {1} <   \dots <   i _ {k}} x _ {i _ {1}} \dots x _ {i _ {k}} = (- 1) ^ {k} a _ {k} \quad (1 \leqslant k \leqslant n).\tag{24}
\]

More precisely, we have

\[
\mathrm{E} _ {f} = \mathrm{A} [ \mathrm{X} _ {1}, \dots , \mathrm{X} _ {n} ] / \mathfrak {a}
\]

where the ideal \(\mathfrak{a}\) is generated by the polynomials \(s_k + (-1)^{k+1} a_k\) for \(1 \leqslant k \leqslant n\) , and \(x_i\) is the residue class of \(\mathbf{X}_i\) mod \(\mathfrak{a}\) for \(1 \leqslant i \leqslant n\) . The relation (24) is also equivalent to \(f(\mathbf{X}) = \prod_{i=1}^{n} (\mathbf{X} - x_i)\) . When there is any risk of ambiguity, we write \(x_{1,f}, \ldots, x_{n,f}\) in place of \(x_1, \ldots, x_n\) .

PROPOSITION 4. --- Let B be a commutative ring, \(\rho\) a homomorphism of A into B and \(\xi_1, \ldots, \xi_n\) elements of B. Suppose that the relation \( {}^{\rho}f(X) = \prod_{i=1}^{n}(X - \xi_i)\) holds in B{[}X{]}. Then there exists one and only one ring homomorphism \(u: E_f \to B\) such that \(\rho(a) = u(a, 1)\) for all \(a \in A\) and \(u(x_i) = \xi_i\) for \(1 \leqslant i \leqslant n\) .

We consider B as associative, commutative and unital A-algebra by means of \(\rho\) . Then the relation \({}^{\rho}f(X) = \prod_{i=1}^{n}(X - \xi_i)\) may also be written in the form

\[
\sum_ {i _ {1} <   \dots <   i _ {k}} \xi_ {i _ {1}} \dots \xi_ {i _ {k}} = (- 1) ^ {k} a _ {k} \cdot 1 \quad (1 \leqslant k \leqslant n)
\]

in B. Since the relations (24) define a presentation of \(\mathbf{E}_f\) , Prop. 4 follows.

Prop. 4 justifies the name « universal decomposition algebra of f» for \(E_{f}\) . The relation \(f(\mathbf{X}) = \prod_{i=1}^{n} (\mathbf{X} - x_{i,f})\) is called the « universal decomposition of f». Let \(\sigma \in S_{n}\) be a permutation; since \(f(\mathbf{X}) = \prod_{i=1}^{n} (\mathbf{X} - x_{\sigma(i),f})\) , there exists an automorphism \(t_{\sigma}\) of the A-algebra \(E_{f}\) characterized by \(t_{\sigma}(x_{i,f}) = x_{\sigma(i),f}\) for \(1 \leqslant i \leqslant n\) . We have \(t_{\sigma\tau} = t_{\sigma} \circ t_{\tau}\) for \(\sigma, \tau\) in \(S_{n}\) , and hence obtain an action of the group \(S_{n}\) on the A-algebra \(E_{f}\) .

PROPOSITION 5. --- In the universal decomposition algebra \(E_{f}\) the family of monomials \(x_{1}^{\nu(1)}\ldots x_{n}^{\nu(n)}\) such that \(0 \leqslant \nu(i) < i\) for \(1 \leqslant i \leqslant n\) is a basis of the A-module \(E_{f}\) . In particular \(E_{f}\) is a free A-module of rank \(n!\) .

Write \(\mathbf{B} = \mathbf{A}[X_1, \ldots, X_n]\) and \(\mathbf{C} = \mathbf{A}[X_1, \ldots, X_n]^{\mathrm{sym}}\) . By Th. 1 (IV, p.~62) we have \(\mathbf{C} = \mathbf{A}[s_1, \ldots, s_n]\) and \(s_1, \ldots, s_n\) are algebraically independent over \(\mathbf{A}\) . The polynomials without constant term in \(s_1, \ldots, s_n\) form an ideal \(\mathbf{C}^+\) in \(\mathbf{C}\) , supplementary to \(\mathbf{A}\) and generated by \(s_1, \ldots, s_n\) . Let \(\mathfrak{c}\) be the ideal of \(\mathbf{C}\) generated by \(s_1 + a_1\) , \(s_2 - a_2\) , \ldots, \(s_n + (-1)^{n+1}a_n\) . There exists an \(\mathbf{A}\) -algebra automorphism of \(\mathbf{C}\) which maps \(s_k\) to \(s_k + (-1)^{k+1}a_k\) for \(1 \leqslant k \leqslant n\) , and hence maps \(\mathbf{C}^+\) onto \(\mathfrak{c}\) ; therefore we have \(\mathbf{C} = \mathbf{A} \oplus \mathfrak{c}\) . Moreover Th. 1, c) of IV, p.~62 shows that

\[
\mathbf {B} = \oplus \mathbf {C X} ^ {\nu}
\]

where S is the set of all \(\nu \in \mathbf{N}^n\) such that \(0 \leqslant \nu(i) < i\) for \(1 \leqslant i \leqslant n\) . The ideal \(\mathfrak{a}\) of B is generated by \(\mathfrak{c}\) , whence \(\mathfrak{a} = \mathbf{B}\mathfrak{c} = \bigoplus_{\nu \in \mathbb{S}} \mathfrak{c} \cdot \mathbf{X}^\nu\) . Since \(\mathbb{C} = \mathbb{A} \oplus \mathfrak{c}\) , we find

\[
\mathbf {B} = \mathfrak {a} \oplus \bigoplus_ {\nu \in S} \mathbf {A X} ^ {\nu},
\]

whence Prop. 5, because \(\mathbf{E}_f = \mathbf{B} / \mathfrak{a}\) .

COROLLARY. --- The canonical homomorphism of A into the universal decomposition algebra of the monic polynomial \(f \in A[X]\) is injective.

For the unit element of \(E_{f}\) forms part of a basis of the A-module \(E_{f}\) .

PROPOSITION 6. --- Let \(f \in \mathbf{A}[\mathbf{X}]\) be a monic polynomial of degree \(n\) , and let \(P\) be a symmetric polynomial in \(\mathbf{X}_1, \ldots, \mathbf{X}_n\) with coefficients in \(A\) . Then there exists precisely one element \(a\) of \(A\) with the following property:

(FS) For any ring homomorphism \(\rho: \mathbf{A} \to \mathbf{B}\) and decomposition \({}^{\rho} f(\mathbf{X}) = \prod_{i=1}^{n} (\mathbf{X} - \xi_i)\) in \(\mathbf{B}[\mathbf{X}]\) we have \(\rho(a) = \mathrm{P}(\xi_1, \ldots, \xi_n)\) .

Write \(f = \mathbf{X}^n + \sum_{k=1}^{n} a_k \mathbf{X}^{n-k}\) , then by Th. 1, IV, p.~62 there exists a polynomial

\(\Pi\) in \(n\) indeterminates with coefficients in \(\mathbf{A}\) , such that \(\mathbf{P} = \Pi(s_1, \ldots, s_n)\) . Put \(a = \Pi(-a_1, a_2, \ldots, (-1)^n a_n)\) . Under the hypothesis (FS) we have

\[
s _ {k} (\xi_ {1}, \dots , \xi_ {n}) = (- 1) ^ {k} \rho (a _ {k})
\]

whence

\[
\begin{array}{r l} \rho (a) & = \Pi (- \rho (a _ {1}), \rho (a _ {2}), \dots , (- 1) ^ {n} \rho (a _ {n})) \\ & = \Pi (s _ {1} (\xi_ {1}, \dots , \xi_ {n}), \dots , s _ {n} (\xi_ {1}, \dots , \xi_ {n})) \\ & = P (\xi_ {1}, \dots , \xi_ {n}). \end{array}
\]

This proves the existence of an element a satisfying (FS). The uniqueness of a follows from the Cor. to Prop. 5, because we have \(a \cdot 1 = P(x_{1,f}, \ldots, x_{n,f})\) in the universal decomposition algebra \(\mathbf{E}_f\) .

With the notation of Prop. 6 one sometimes writes \(a = \mathrm{P}^{\#}(f)\) . Here are some examples.

Examples. --- 1) If \(\mathrm{P} = s_k\) , then \(\mathrm{P}^{\#}(f) = (-1)^k a_k\) .

2) Let \(g\) be a polynomial in \(\mathbf{A}[\mathbf{X}]\) and put

\[
\mathrm{P} \left(\mathrm{X} _ {1}, \dots , \mathrm{X} _ {n}\right) = g \left(\mathrm{X} _ {1}\right) \dots g \left(\mathrm{X} _ {n}\right).
\]

Then \(\mathrm{P}^{\#}(f)\) is just the resultant \(\operatorname{res}(f, g)\) , by Cor. 1 of IV, p.~80.

\begin{enumerate}
\def\labelenumi{\arabic{enumi})}
\setcounter{enumi}{2}
\item
  Put \(\Delta(X_1, \ldots, X_n) = \prod_{i < j} (X_i - X_j)^2\) , then \(\Delta^{\#}(f)\) is just the discriminant of the monic polynomial \(f\) (IV, p.~82, formula (46)).
\item
  Put \(\mathrm{P}(\mathrm{X}_{1},\ldots,\mathrm{X}_{n})=\mathrm{X}_{1}^{k}+\cdots+\mathrm{X}_{n}^{k}\) ; moreover, define the algebra \(\mathrm{E}=\mathrm{A}[\mathrm{X}]/(f)\) and write x for the image of X in E. We recall that the A-module E is free, with basis \((1,x,\ldots,x^{n-1})\) (IV, p.~11, Cor.). Let us show
\end{enumerate}

\[
\operatorname{Tr} _ {\mathrm{E/A}} (x ^ {k}) = \mathrm{P} ^ {\#} (f).\tag{25}
\]

Put \(\pi_k = \mathrm{Tr}_{\mathrm{E / A}}(x^k)\) for every integer \(k \geqslant 1\) . Bearing in mind Newton's relations (IV, p.~70) we need only establish the relations

\begin{enumerate}
\def\labelenumi{(\arabic{enumi})}
\setcounter{enumi}{25}
\tightlist
\item
\end{enumerate}

\[
\begin{array}{l l} \pi_ {k} + a _ {1} \pi_ {k - 1} + \dots + a _ {k - 1} \pi_ {1} + k a _ {k} = 0 & \text { for } \quad 1 \leqslant k \leqslant n \\ \pi_ {k} + a _ {1} \pi_ {k - 1} + \dots + a _ {n - 1} \pi_ {k - n + 1} + a _ {n} \pi_ {k - n} = 0 & \text { for } \quad k > n \end{array}\tag{27}
\]

(which we shall also call « Newton's relations »). The relation (27) is clear, because the left-hand side is the trace of

\[
x ^ {k - n} \left(x ^ {n} + a _ {1} x ^ {n - 1} + \dots + a _ {n - 1} x + a _ {n}\right) = 0.
\]

Suppose that \(1 \leqslant k \leqslant n\) and put

\[
y = x ^ {k} + a _ {1} x ^ {k - 1} + \dots + a _ {k - 1} x + a _ {k} \cdot 1;
\]

let \(\mathbf{M} = (m_{ij})\) be the matrix of the linear mapping \(u\mapsto yu\) in E relative to the basis \((x^{i})_{0\leqslant i\leqslant n-1}\) . We easily obtain the relations

\[
\begin{array}{l l} m _ {i i} = a _ {k} & \text { for } \quad 0 \leqslant i <   n - k \\ m _ {i i} = 0 & \text { for } \quad n - k \leqslant i <   n, \end{array}
\]

whence

\[
\operatorname{Tr} _ {\mathrm{E/A}} (y) = \sum_ {i = 0} ^ {n - 1} m _ {i i} = (n - k) a _ {k}.
\]

Moreover we have

\[
\operatorname{Tr} _ {\mathrm{E/A}} (y) = \pi_ {k} + a _ {1} \pi_ {k - 1} + \dots + a _ {k - 1} \pi_ {1} + n a _ {k},
\]

whence formula (26) follows.

